\section{Related work}

\subsection{Articulation work and common information spaces}

Articulation work concerns the additional coordination required to make cooperative activity possible when plans, actors, resources, and contingencies do not align automatically \cite{schmidt1992articulation}. In connected laboratories, this work includes eliciting omitted constraints, reconciling different vocabularies, identifying which representation governs execution, and determining who can act on a device, method, or sample. The work is distributed among users, maintainers, vendors, integrators, and scientific operators. Public discussions provide a partial common information space, although access to relevant manuals, licences, configurations, and private outcomes remains uneven.

Technical question-and-answer communities have been studied as environments for problem solving and knowledge sharing \cite{treude2011questions,vasilescu2014social}. The present setting differs from ordinary programming support because evidence is tied to instruments, material state, physical geometry, calibration, and scientific consequence. Participants must coordinate what software recorded with what physically occurred.

\subsection{Infrastructure and breakdown}

Infrastructure is relational, embedded in practice, and frequently becomes visible during breakdown \cite{star1996infrastructure}. Laboratory automation illustrates this property across device protocols, vendor databases, resource definitions, training, support, and local conventions. A failed upgrade reveals the dependency on a legacy interface. A version-control incident reveals that an apparent source directory contains active recovery state. A cross-machine tip failure reveals hidden coordinate semantics and local teaching. The forum makes these dependencies visible by assembling evidence around the breakdown.

Infrastructuring also concerns the ongoing work through which communities build and maintain shared resources. A useful answer can remain confined to one thread, move to a restricted channel, or become a maintained artifact. The latter transition requires work on representation, ownership, versioning, validation, and correction. We examine this transition as a central outcome of public technical collaboration.

\subsection{Repair and material consequence}

Repair sustains sociotechnical systems by responding to breakdown and re-establishing workable order \cite{jackson2014repair}. In laboratory automation, repair can involve diagnosis, manual intervention, software recovery, state reconciliation, material salvage, and a judgment about scientific validity. The relevant object is not only the software process. It includes samples, liquid in tips, device positions, contamination histories, and timing constraints.

This materiality complicates rollback. Once aspiration, dispensing, motion, heating, or incubation occurs, a software restart cannot restore the prior physical world. Repair therefore depends on evidence about the completed action prefix and the uncertainty that remains. The public record also distinguishes repair of the immediate run from repair of the infrastructure: a local workaround may restore operation without producing a durable public correction.

\subsection{Knowledge stewardship and open technical communities}

Open-science communities support peer learning, access, and shared practice \cite{eerland2021communities}. FAIR principles emphasize findability, accessibility, interoperability, and reusability \cite{wilkinson2016fair}. Public technical communities add a governance problem. Knowledge can become stale, version-specific, disputed, privately migrated, or detached from its evidence. Durable stewardship requires provenance, applicability, validation stage, maintenance ownership, last verification, known limitations, and supersession.

Our analysis connects these concerns. We study how participants produce the context required for cooperative diagnosis, how repair crosses software and material state, and how public knowledge becomes or fails to become maintained infrastructure.
